\scanpage{83}
% archive foliation: 82
\deleted{$\unsure{X_0/S}$} d'une situation $X_0/I_0/\unsure{x_0|k_0}/k_0$
où $k_0$ est un sous-corps. [OPS \unsure{disons} que
$k$ est une clôture alg. de $k_0\ldots$].
Alors $\Gamma=\operatorname{Gal}(\overset{\unsure{\longleftarrow|\longleftrightarrow}}{k/k_0})$ opère sur
$\Sigma^{\mathrm{ext}}$ par transport de structure, donc on a
une \deleted{\unsure{la}} opération aussi sur $P(X/I/k)$,
\deleted{donc}\inserted{ou encore} un hom
\[
\Gamma\longrightarrow{}^{P(X/I/k)}\widetilde\Gamma_\nu
\simeq\operatorname{Aut}(\Sigma^{\mathrm{ext}})
\]
qui par image inverse \inserted{d'extensions} donne naissance
à l'extension
\[
1-\pi_1(X\setminus I,x)-\pi_1(X_0\setminus I_0,x_0)
\longrightarrow\pi_1(k_0/k)-1
\]
$\ldots\quad\Gamma\xrightarrow{\chi}\operatorname{Aut}(T)\simeq\widehat{\mathbb Z}^{*}$
est l'hom de Kummer-Tate $\ldots$

\textcircled{II} \uline{Plus} précisément supposons $k_0=\mathbb Q$, $X_0=\mathbb P^1_{\mathbb Q}$,
\inserted{$\overline{\mathbb Q}$ clôture alg. de $\mathbb Q$ dans $\mathbb C$},
\[
I=\mu_{\nu+1}(\overline{\mathbb Q})=\mu_{\nu+1}(\mathbb C^{*})
\subset\overline{\mathbb Q}
=\unsure{\mathbb E|\mathbb A}^1_{\overline{\mathbb Q}}(\overline{\mathbb Q})
\subset\mathbb P^1_{\overline{\mathbb Q}}(\overline{\mathbb Q})
\]
On trouve un isomorphisme canonique,
via l'exponentielle et les lacets
\deleted{$\unsure{I\simeq[0,\nu]|I\simeq[0,1]}$},
\[
\Sigma(X/I/k)^{\mathrm{ext}}\simeq L_\nu^{\mathrm{ext}}
\]
donc $P(X,I/k)=$ torseur trivial
